\section*{导言}

这场访谈的特别之处，在于它几乎把一个顶尖研究者近十多年的问题演化完整摊开了看。从交大 ACM 班、本科实习、博士训练，到 FAIR、自监督、DiT、Cambrian、JEPA 与 AMI Labs，谢赛宁并不是在讲一串彼此孤立的成就，而是在不断回答同一个问题：如果我们真想理解智能，应该把注意力放在哪些变量上？因此，这份笔记不按时间顺序复述，而是按十个主题重新组织，希望把访谈中那些跨越不同年份、却在逻辑上彼此呼应的判断压缩成更清晰的结构。

如果只把这场访谈当成“明星科学家访谈”，会错过它最珍贵的部分。它真正罕见的地方，是谢赛宁愿意把很多本可被包装得更漂亮的经历，重新还原成充满偶然、失败、误判、迟到与反复修正的过程。他讲导师、论文、创业、世界模型时，始终在反复拆解一个误区：外界看到的是结果的光滑表面，而研究者真正需要面对的是形成结果之前那团极不规则的混沌。这也是为什么他一直强调非线性、探索、负面结果与反脆弱，因为这些概念恰好构成了从混沌里长出新结构的必要条件。

访谈的第二层价值，在于它让“视觉研究”重新获得了思想上的重量。在大模型语境里，很多人会不自觉地把 vision 理解为语言模型的外围组件，或者把 world model 理解为一轮新的叙事包装。但谢赛宁给出的论证非常明确：视觉之所以重要，不是因为它能给 LLM 再加一个模态，而是因为它连接着真实世界、连续空间、时间流和行动后果。换句话说，若把智能只理解成在 token 空间里生成下一个符号，那么整个 AI 领域就很容易把“会说”误当成“会理解”。

\begin{importantbox}{阅读这份笔记的三条主线}
第一条主线是“人成为研究者的过程”，包括家庭、学校、导师和同辈影响如何共同塑形。第二条主线是“视觉通向世界模型的理论连续性”，它解释了为什么谢赛宁从未把自己的工作看作离散跳跃。第三条主线是“研究与组织如何相互支撑”，从 FAIR 到 AMI Labs，真正变化的不是问题本身，而是承载问题的制度环境。
\end{importantbox}

最后，这份笔记有意保留访谈中的文化与人生部分。电影、哲学、纽约、动物智能、维特根斯坦、克洛普、书单与命运感，并不是技术讨论的装饰物，而是谢赛宁理解智能时不可分割的背景。若没有这些背景，他对 research taste、对语言边界的警惕、对真实世界的坚持，就会被误读成几句漂亮口号。正因为他始终把技术放回更大的生活世界中衡量，这场访谈才不只是在谈 AI，也是在谈怎样在一个高速竞争、叙事过热的时代，尽量保持清醒。

